\documentclass[11pt, letterpaper]{article}
\usepackage[utf8]{inputenc}
\usepackage[pdftex]{graphicx}
\usepackage{epstopdf}
\usepackage{import}
\usepackage{array}
\usepackage{float}
\graphicspath{{Images/}}

\usepackage{amssymb}    % http://ctan.org/pkg/amssymb
\usepackage{pifont}     % http://ctan.org/pkg/pifont

\usepackage{multirow}
\usepackage{booktabs}
\setlength{\parindent}{1.5em}
\setlength{\tabcolsep}{2pt}
\usepackage[export]{adjustbox}

\usepackage[english, spanish]{babel}
\usepackage[nottoc, numbib]{tocbibind}
\renewcommand{\rmdefault}{phv} % Arial
\renewcommand{\sfdefault}{phv} % Arial
\usepackage[left=3cm,right=2.5cm,top=3cm,bottom=2.5cm]{geometry}

\title{Diseño y Desarrollo de una Cama de Prueba (TestBed) de una Red Definida por Software para el Transporte de Paquetes de una red IMS con QOS}
\author{
    Jonathan Dannilo Arcila Muñoz
\and
	Wilson Bernardo Benavides Cepeda    
}
\date{2017}

%========================START DOCUMENT========================%

\begin{document}
\selectlanguage{spanish}

%-------------------------TITLE PAGE---------------------------%

\begin{titlepage}
\thispagestyle{empty}
\centering

\LARGE{\textbf{Diseño y Desarrollo de una Cama de Prueba (TestBed) de una Red Definida por Software para el Transporte de Paquetes de una red IMS con QOS}}\\
\vspace{1.8cm}

\includegraphics[scale=1]{logo-unicauca.png}

\begin{center}
\Large
\vspace{8mm}
Anteproyecto de Trabajo de Grado\\
\vspace{1.2cm}
\textbf{Jonathan Dannilo Arcila Muñoz \\
Wilson Bernardo Benavides Cepeda}\\
\vspace{1.2cm}
Director: Ph.D. Oscar Mauricio Caicedo Rendón\\


\vfill
Universidad del Cauca\\
Facultad de Ingeniería Electrónica y Telecomunicaciones\\
Departamento de Telemática\\
Popayán, Junio, 2018

\end{center}
\end{titlepage}
%-------------------TABLE OF CONTENTS (Roman Numeration) ------------------------%

\pagestyle{plain}
\pagenumbering{roman}
\setcounter{page}{3}
\tableofcontents
\break

%--------------------- DOCUMENT BODY ------------------------%
\pagenumbering{arabic}

\section{Planteamiento del Problema}
	\import{Problem/}{problem_statement.tex}

%--------------------- BIBLIOGRAPHY ------------------------%	
\bibliographystyle{IEEEtran}
%\bibliographystyle{plain}
\bibliography{bibliography}

\end{document}






















